\documentclass[a4paper,11pt]{article}
\usepackage{exam-style-842-unit}
\examinehead{842 信息化概论 · 第 4 单元练习册}{练习一 参考答案与解析}
\begin{document}

\unitanswerscover{842 信息化概论 · 单元练习册}{第 4 单元参考答案}{练习一（软件工程与开发流程）的答案与解析；每道判断题均给出改正后的正确表述。教材页码按“扫描页−7＝印刷页”换算，引文以 OCR 原文为准。}

\answerpart{一、单项选择题（10 题，每题 3 分，共 30 分）}

\q B。\begin{analysis}1968 年北大西洋公约组织（NATO）的计算机科学家在联邦德国召开的国际学术会议上第一次提出“软件危机”这个名词，同时第一次提出“软件工程”概念（教材 3.3.1，第 39 页）。A 国际标准化组织（ISO）在软件工程部分给出的是 ISO 9000 的定义；C IEEE 给出的是软件工程术语汇编中的定义；D 世界电子商务会议提出的是电子商务概念（教材 8.1.1）。\end{analysis}

\q B。\begin{analysis}教材 3.3.1 把软件开发分为 3 个阶段：1950—1960 年为程序设计阶段，1960—1970 年为软件设计阶段，1970 年至今为软件工程时代。其中 1960—1970 年采用小组合作生产方式，出现“软件作坊”，软件开始作为产品被广泛使用（第 38 页）。A、C 时间对应错误；D 面向对象属于软件工程时代采用的技术方法，不是教材的阶段划分。\end{analysis}

\q B。\begin{analysis}软件过程（Software Process）是指为建造高质量软件所需完成的任务的框架，即形成软件产品的一系列步骤，包括中间产品、资源、角色及过程中采取的方法、工具等范畴（教材 3.3.3，第 40 页）。A 软件工程是工程学科与总体方法论；C 软件项目管理是组织管理活动；D 软件体系结构关注系统的组成结构，三者都不是“任务框架”的定义。\end{analysis}

\q B。\begin{analysis}1970 年温斯顿·罗伊斯（Winston Royce）提出了著名的“瀑布模型”，直到 20 世纪 80 年代早期，它一直是唯一被广泛采用的软件开发模型（教材 3.3.3，第 40 页）。A 布鲁克斯是《人月神话》的作者；C 迈克尔·波特提出价值链理论（教材 8.1.1）；D 冯·诺依曼提出的是计算机体系结构思想。\end{analysis}

\q A。\begin{analysis}瀑布模型将软件生命周期划分为需求分析、架构设计、详细设计、编码、测试和维护 6 个基本活动，并规定它们自上而下、相互衔接的固定次序，如同瀑布流水逐级下落（教材 3.3.3，第 40 页）。B 把详细设计置于架构设计之前；C 架构设计先于需求分析；D 编码先于详细设计，均与原文顺序不符。\end{analysis}

\q A。\begin{analysis}《人月神话》是布鲁克斯在总结 IBM360 系列机操作系统开发教训的基础上写成的软件工程经典著作，书中对软件开发过程的 5 个重要关键点提出了独到见解（教材 3.3.1，第 39 页）。教材后文虽以 Windows 操作系统 1500 万行以上代码说明软件系统的复杂性（3.3.4），但它不是《人月神话》总结的对象；UNIX、阿波罗登月软件也不符合教材表述。\end{analysis}

\q A。\begin{analysis}软件项目管理是为了使软件项目能够按照预定的成本、进度、质量顺利完成，而对人员（People）、产品（Product）、过程（Process）和项目（Project）进行分析和管理的活动（教材 3.3.4，第 41 页）。B、C、D 均不是教材给出的 4P。\end{analysis}

\q C。\begin{analysis}20 世纪 70 年代中期，美国国防部专门研究软件开发不能按时提交、预算超支和质量达不到用户要求的原因，结果发现 70\% 的项目是因为管理不善引起的，而非技术原因（教材 3.3.4，第 41 页）。可作为补充记诵的还有：到 20 世纪 90 年代中期，大约只有 10\% 的项目能够在预定费用和进度下交付；1995 年美国取消了 810 亿美元的商业软件项目，其中 31\% 未做完就被取消，53\% 的进度延长 50\% 的时间，只有 9\% 能及时交付且费用控制在预算内。\end{analysis}

\q A。\begin{analysis}CMM 共 5 个等级：初始级、可重复级、已定义级、定量管理级、优化级。初始级的软件过程是无序的，有时甚至混乱，对过程几乎没有定义，成功只取决于个人努力，管理是反应式的（教材 3.3.4，第 42 页）。B 可重复级已建立基本的项目管理过程；C 已定义级已实现过程文档化、标准化；D 优化级是最高等级。\end{analysis}

\q D。\begin{analysis}软件项目管理的内容主要包括：人员的组织与管理、软件度量、软件项目计划、风险管理、软件质量保证、软件过程能力评估、软件配置管理（教材 3.3.4，第 41 页）。计算机硬件采购属于硬件与采购管理范畴，不是软件项目管理的内容，故 D 符合题意。\end{analysis}

\answerpart{二、多项选择题（4 题，每题 4 分，共 16 分）}

\q ABCD。\begin{analysis}《人月神话》对软件开发过程提出的 5 个关键点是：（1）提倡外科手术式的团队组织；（2）软件项目的核心概念要由很少的人来完成，以保证概念的完整性；（3）软件开发过程中必要的沟通手段；（4）如何保持适度的文档；（5）在软件开发的过程中要适度改进（教材 3.3.1，第 39 页）。A、B、C、D 分别对应第（1）（2）（3）（4）点，四项全对。\end{analysis}

\q ABC。\begin{analysis}瀑布模型的优点有 4 条：（1）通过设置里程碑，明确每阶段的任务与目标；（2）可为每阶段制订开发计划，进行成本预算，组织开发力量；（3）通过阶段评审，将开发过程纳入正确轨道；（4）严格的计划性保证软件产品的按时交付（教材 3.3.3，第 40—41 页）。D“缺乏灵活性，不能适应用户需求的改变”是瀑布模型的缺点而非优点，故排除。\end{analysis}

\q ABCD。\begin{analysis}瀑布模型的缺点有 4 条：（1）缺乏灵活性，不能适应用户需求的改变；（2）开始阶段的小错误被逐级放大，可能导致软件产品报废；（3）返回上一级的开发需要十分高昂的代价；（4）随着软件规模和复杂性的增加，软件产品成功的机率大幅下降（教材 3.3.3，第 41 页）。四项全对。教材同时指出，尽管受到很多批评，瀑布模型对很多类型的项目依然有效，正确使用可节省大量时间和金钱。\end{analysis}

\q ABC。\begin{analysis}A 正确，CMM 的全称是能力成熟度模型（Capability Maturity Model for Software，SW-CMM）；B 正确，CMM 的核心是把软件开发视为一个过程，并据此对软件开发和维护进行过程监控与研究；C 正确，CMM 分为初始级、可重复级、已定义级、定量管理级、优化级 5 个等级，优化级最高（教材 3.3.4，第 41—42 页）。D 错误：CMM 是对“软件组织”在定义、实施、度量、控制和改善其软件过程实践中各个发展阶段的描述，评价的是组织的过程能力，而非个人程序员的编程能力。\end{analysis}

\answerpart{三、判断题（6 题，每题 2 分，共 12 分）}

\q ×。改正：软件工程是将系统化、严格约束的、可量化的方法应用于软件的开发、运行和维护的工程学科，覆盖需求分析、设计、编码、测试、维护的软件全生命周期，并非只关注编码环节。\begin{analysis}教材 3.3.2 给出的通行定义是：软件工程是研究和应用如何以系统性的、规范化的、可定量的过程化方法去开发和维护软件，以及如何把经过时间考验而证明正确的管理技术和当前能够得到的最好的技术方法结合起来（第 40 页）。把软件工程等同于编程，是软件开发早期“程序设计阶段”的观念。\end{analysis}

\q √。\begin{analysis}瀑布模型的核心思想是按工序将问题化简，将功能的实现与设计分开，把软件生命周期划分为需求分析、架构设计、详细设计、编码、测试和维护 6 个基本活动，并规定它们自上而下、相互衔接的固定次序，如同瀑布流水逐级下落（教材 3.3.3，第 40 页）。\end{analysis}

\q ×。改正：“软件危机”和“软件工程”是 1968 年北约（NATO）计算机科学家会议第一次提出的；1970 年温斯顿·罗伊斯提出的是瀑布模型。\begin{analysis}教材 3.3.1 记载：1968 年北约组织的计算机科学家在联邦德国召开的国际学术会议上第一次提出“软件危机”这个名词，并第一次提出“软件工程”概念；教材 3.3.3 记载：1970 年罗伊斯提出瀑布模型。两个事件的时间与提出者不能混淆。\end{analysis}

\q ×。改正：软件项目不能靠简单增加人手来解决进度问题；《人月神话》主张外科手术式的团队组织，核心概念要由很少的人完成以保证概念的完整性，进度落后时盲目加人往往使项目更加落后。\begin{analysis}教材 3.3.1 指出：软件开发组织上的过分民主往往带来没有效率和责任，应该借鉴外科手术式的团队方式，有一个主要责任人，其他人分工协作；核心概念要少而精，太多想法会使项目没有焦点（第 39 页）。这正对应“人月神话”所揭示的“加人不能无限缩短工期”。\end{analysis}

\q √。\begin{analysis}教材 3.3.4 指出：软件项目管理的根本目的是让软件项目尤其是大型项目的整个软件生命周期（从分析、设计、编码到测试、维护全过程）都能在管理者的控制之下，以预定成本按期、按质地完成软件交付用户使用（第 41 页）。\end{analysis}

\q √。\begin{analysis}CMM 的 5 个等级中，优化级是最高等级，其过程的量化反馈和先进的新思想、新技术促使过程不断改进（教材 3.3.4，第 42 页）。教材同时指出 CMM 在具体管理上过于繁琐和复杂，且近年有弱化的趋势，但这不影响“优化级为最高等级”的表述。\end{analysis}

\answerpart{四、配对题（6 对，每对 1 分，共 6 分）}

\q （1）—A\quad （2）—B\quad （3）—C\quad （4）—D\quad （5）—E\quad （6）—F。
\begin{analysis}1968 年北约会议首次提出“软件危机”与“软件工程”；1970 年罗伊斯提出瀑布模型；布鲁克斯总结 IBM360 操作系统开发的教训并著《人月神话》；SW-CMM 即能力成熟度模型，分 5 个等级；软件过程是形成软件产品的一系列步骤（任务框架）；软件配置管理是软件项目管理的内容之一（教材 3.3.1、3.3.3、3.3.4）。\end{analysis}

\answerpart{五、简答题（3 题，每题 8 分，共 24 分）}

\q 参考答案要点：
（1）软件工程概念（2 分）：软件工程一直以来缺乏统一定义，教材列举了多种代表性定义，一般认为软件工程是研究和应用如何以系统性的、规范化的、可定量的过程化方法去开发和维护软件，并把经过时间考验而证明正确的管理技术与当前最好的技术方法结合起来。
（2）两种代表性定义（4 分）：其一，IEEE 在软件工程术语汇编中的定义——软件工程是将系统化的、严格约束的、可量化的方法应用于软件的开发、运行和维护，即将工程化应用于软件；其二，《计算机科学技术百科全书》的定义——软件工程是应用计算机科学、数学、逻辑学及管理科学等原理开发软件的工程，借鉴传统工程的原则、方法，以提高质量、降低成本和改进算法为目的。另可补充 ISO 9000 定义：软件工程是输入转化为输出的一组彼此相关的资源和活动。
（3）两个重要问题（2 分）：一般认为，软件工程中重要的问题是如何解决软件过程和进行软件管理；即一手抓“软件过程”（任务框架与阶段活动），一手抓“软件项目管理”（人员、产品、过程、项目 4P）。（教材 3.3.2，第 40 页）

\q 参考答案要点：
（1）软件危机的含义（2 分）：软件开发中出现的开发时间大大超出规划时间表、预算超支、质量达不到用户要求、项目大量失败等尖锐矛盾，即软件危机。
（2）主要表现（2 分）：进度难以控制、成本居高不下、质量难以保证、维护困难、失败项目屡见不鲜。教材以 IBM360 操作系统为最突出的例子，项目负责人布鲁克斯形容程序设计工作像“泥潭”。
（3）产生原因（2 分）：软件规模急剧膨胀、复杂性增加，而软件开发长期沿用个体化、手工作坊式的开发方式，开发技术没有突破性进展，开发人员的开发技术不适应规模大、结构复杂的软件开发；同时软件是纯知识产品，开发进度和质量难以估计和度量。
（4）产生背景（2 分）：1968 年北约计算机科学家会议第一次提出“软件危机”和“软件工程”，软件工程学作为研究和克服软件危机的新兴工程学科应运而生；此后软件开发进入工程化生产时代。（教材 3.3.1，第 38—39 页）

\q 参考答案要点：
（1）核心思想（2 分）：按工序将问题化简，将功能的实现与设计分开；把软件生命周期划分为 6 个基本活动，并规定自上而下、相互衔接的固定次序，如同瀑布流水逐级下落。
（2）6 个基本活动（2 分）：需求分析$\rightarrow$架构设计$\rightarrow$详细设计$\rightarrow$编码$\rightarrow$测试$\rightarrow$维护。
（3）优点（2 分）：通过设置里程碑明确每阶段任务与目标；可为每阶段制订开发计划、进行成本预算、组织开发力量；通过阶段评审把开发过程纳入正确轨道；严格的计划性保证软件产品按时交付。
（4）缺点（2 分）：缺乏灵活性，不能适应用户需求改变；开始阶段的小错误被逐级放大，可能导致产品报废；返回上一级开发需付出十分高昂的代价；随着软件规模和复杂性增加，成功机率大幅下降。（教材 3.3.3，第 40—41 页）

\answerpart{六、综合题（1 题，共 12 分）}

\q 参考答案要点：
（1）问题诊断（6 分，每点 1 分，答出 5 点即可）：
第一，盲目加人无效，违背《人月神话》的规律——软件开发组织应借鉴外科手术式团队，核心概念由很少的人完成以保证概念完整性，简单增加人手反而增加沟通成本、降低效率。
第二，需求频繁变更导致返工，说明选用瀑布模型不当——瀑布模型缺乏灵活性，不能适应用户需求的改变，且返回上一级开发代价高昂、早期小错误会被逐级放大。
第三，进度失控说明缺少里程碑、阶段评审与严格的开发计划——瀑布模型的优点正是通过里程碑明确阶段目标、通过阶段评审把过程纳入正轨，该项目未落实。
第四，缺少必要的沟通手段——《人月神话》指出软件开发中最大的风险往往不是技术缺陷，而是缺少沟通。
第五，文档与概念完整性管理不足——应保持适度文档（文档不是最终产品，但过多或过少都不可取），并通过配置管理控制需求与版本变更。
第六，管理不善是主因——教材 3.3.4 指出 70\% 的失败项目源于管理不善而非技术原因，该项目未开展风险管理、软件度量与质量保证。
（2）改进建议（6 分，每点 1.5 分，答出 4 点即可）：
第一，调整团队组织，设一名主要责任人统筹，核心业务概念由少数骨干确定，减少无效沟通。
第二，控制需求变更：建立基线，变更走评审与配置管理流程；需求不确定时改用原型法或迭代式开发，而非硬套瀑布模型。
第三，设置里程碑与阶段评审，制订分阶段开发计划与成本预算，及时纠偏。
第四，引入风险管理、软件度量与质量保证，用数据监控进度和质量，并参照 CMM 逐步提升过程能力。
第五，保持适度的文档与代码规范，做好配置管理与版本控制。
（3）结论（1 分）：软件项目既要选对开发过程模型（如需求稳定的项目可用瀑布模型，需求多变宜用迭代/原型），又要落实项目管理与团队组织原则；进度落后时先诊断管理问题，而不是简单增加人手。
\end{document}
